\chapter{The Fourier Transform}

In the last chapter we learned to take a \emph{periodic} signal and break it into a
discrete ladder of harmonics: a sum of sinusoids whose frequencies are integer
multiples of one fundamental. That is the Fourier \emph{series}. But almost nothing
in the real world is exactly periodic. A hand clap, a footstep, a single twitch of a
muscle, the chirp of a bird, the pulse you send down a wire -- these are
\emph{transient} events. They happen once and are gone. They have no period at all.

The Fourier transform is what the Fourier series becomes when we let the period grow
without bound. The discrete ladder of harmonics fuses into a continuous \emph{spectrum}.
Instead of asking ``how much of harmonic number $7$ is in this signal?'' we ask ``how
much of \emph{every} frequency, continuously, is in this signal?'' The answer is itself
a function -- a function of frequency rather than of time. Learning to read that
function, and to see a time-domain wiggle and its frequency-domain description as two
views of one object, is a central skill in signal processing.

This chapter builds the transform geometrically. We will watch the harmonics crowd
together as the period stretches, define the continuous-time Fourier transform (CTFT)
and its inverse, learn to read magnitude and phase, compute the handful of transforms
that you will reach for again and again (especially the rectangle and its sinc), meet
the uncertainty principle that ties time-width to frequency-width, and close the loop
with Parseval's energy theorem. Throughout, the through-line of the book -- a
wrist-worn gesture controller that emits MIDI -- gives us a reason to care: a gesture
is a transient, and the transform is how we will look at it.

\section{From series to transform: letting the period grow without bound}
\label{ch07:series-to-transform}

Recall the complex (exponential) form of the Fourier series for a signal $x(t)$ that
repeats every $T$ seconds. With fundamental angular frequency $\omega_0 = 2\pi/T$,
\begin{equation}
\label{ch07:fs-pair}
x(t) = \sum_{k=-\infty}^{\infty} c_k\, e^{\jj k \omega_0 t},
\qquad
c_k = \frac{1}{T}\int_{-T/2}^{T/2} x(t)\, e^{-\jj k \omega_0 t}\, \dd t .
\end{equation}
Each coefficient $c_k$ is a complex number telling us the amplitude and phase of the
harmonic sitting at frequency $k\omega_0$. The harmonics live on a ladder whose rungs
are spaced $\omega_0 = 2\pi/T$ apart. A short period $T$ means widely spaced rungs; a
long period means a fine ladder.

Now imagine a signal that is really just a single isolated pulse -- say a smooth bump
of muscle activity lasting a fraction of a second -- but which we \emph{pretend} repeats
with some huge period $T$. As we slide $T$ to larger and larger values, two things
happen at once.

\begin{itemize}
  \item The rungs of the ladder, spaced $\omega_0 = 2\pi/T$, move closer together.
        In the limit $T\to\infty$ they merge into a \emph{continuum} of frequencies.
  \item Each coefficient $c_k$, which carries a factor $1/T$, shrinks toward zero. The
        signal is spread out over a longer and longer period, so any single harmonic
        carries vanishingly little of it.
\end{itemize}

The trick is to track the right quantity so that nothing collapses to zero. Define
\begin{equation}
\label{ch07:envelope}
X(\jj\omega) \;=\; T\,c_k \Big|_{\,k\omega_0 = \omega}
\;=\; \int_{-T/2}^{T/2} x(t)\, e^{-\jj \omega t}\, \dd t .
\end{equation}
We multiplied $c_k$ by $T$ to kill the troublesome $1/T$, and we relabeled the discrete
index $k\omega_0$ as a continuous variable $\omega$. The quantity $X(\jj\omega)$ is the
\emph{envelope} that the scaled coefficients trace out. It does not collapse as
$T\to\infty$; it converges to a perfectly good function of $\omega$.

\begin{intuitionbox}
Think of the spectrum as a row of vertical bars, one per harmonic, with height
$\abs{c_k}$. As $T$ grows, the bars get shorter (height $\sim 1/T$) but also more
numerous and densely packed. If instead you plot $T\abs{c_k} = \abs{X(\jj\omega)}$,
the bars settle onto a fixed \emph{curve}. The Fourier transform is that curve: the
smooth silhouette the harmonics approach as the period grows.
\end{intuitionbox}

To recover $x(t)$ we take the synthesis sum \eqref{ch07:fs-pair} and substitute
$c_k = X(\jj k\omega_0)/T = X(\jj k\omega_0)\,\omega_0/(2\pi)$:
\begin{equation}
x(t) = \sum_{k=-\infty}^{\infty} \frac{\omega_0}{2\pi}\, X(\jj k\omega_0)\, e^{\jj k\omega_0 t}.
\end{equation}
This is a Riemann sum. The samples are spaced $\Delta\omega = \omega_0$ apart, the
summand is $\frac{1}{2\pi} X(\jj\omega)e^{\jj\omega t}\,\Delta\omega$, and as
$\omega_0\to 0$ the sum becomes an integral:
\begin{equation}
\label{ch07:ift}
x(t) = \frac{1}{2\pi}\int_{-\infty}^{\infty} X(\jj\omega)\, e^{\jj\omega t}\, \dd\omega .
\end{equation}

\begin{figure}[H]
\centering
\begin{tikzpicture}
\begin{axis}[width=12.5cm,height=5.4cm,
  xlabel={frequency $\omega$}, ylabel={$T\,|c_k|$},
  xmin=-6.3,xmax=6.3,ymin=0,ymax=2.2,
  ytick=\empty, xtick={-6.28,-3.14,0,3.14,6.28},
  xticklabels={$-2\pi$,$-\pi$,$0$,$\pi$,$2\pi$},
  axis lines=left, tick align=outside,
  title style={font=\small}, title={Harmonic bars settling onto the transform envelope $|X(\jj\omega)|$}]
\addplot[specred,thick,domain=-6.28:6.28,samples=181] {abs(2*sin(deg(x/2+0.0001))/(x+0.0001))};
\addplot+[ycomb,sigblue,thick,mark=*,mark size=1.3pt,
  samples at={-6.28,-4.71,-3.14,-1.57,0.0001,1.57,3.14,4.71,6.28}]
   {abs(2*sin(deg(x/2+0.0001))/(x+0.0001))};
\end{axis}
\end{tikzpicture}
\caption{As the assumed period $T$ grows, the harmonics (blue stems) crowd together and
the scaled coefficients $T|c_k|$ trace the continuous envelope $|X(\jj\omega)|$ (red).
The Fourier transform is the limiting silhouette of the line spectrum.}
\label{ch07:fig-envelope}
\end{figure}

\begin{historybox}
The step from the discrete harmonic ladder of a Fourier series to the continuous integral transform
took shape in the early nineteenth century through the work of \textbf{Augustin-Louis Cauchy}, \textbf{Sim\'{e}on Denis
Poisson}, and \textbf{Joseph Fourier}. In his 1822 treatise \emph{Th\'{e}orie analytique de la
chaleur}, Fourier considered heat spreading not in a closed ring of finite length, but along an
infinitely long metal bar extending from $-\infty$ to $+\infty$. By taking the limit as the
period $T \to \infty$, the harmonic spacing $\Delta\omega = 2\pi/T$ shrank to the infinitesimal
$\dd\omega$, and the discrete sum collapsed into the continuous Fourier integral.

Independently in 1815, \textbf{Cauchy}, working on the propagation of surface waves on liquids, had
found reciprocal integral formulas linking spatial wave patterns to their wavenumber
densities. For decades, however, mathematicians disputed the validity of the transform for
non-decaying functions. A pure, eternal sinusoid $\cos(\omega_0 t)$ has infinite total energy and
fails to decay as $\abs{t} \to \infty$; its classical Riemann integral $\int_{-\infty}^\infty
\cos(\omega_0 t) e^{-\jj\omega t}\,\dd t$ diverges. Traditional analysts dismissed transforms
of sinusoids and steps as mathematical absurdities.

The theoretical foundation was finally secured in 1910 by the Swiss mathematician \textbf{Michel
Plancherel}, who proved that the Fourier transform is a unitary, energy-preserving operator on the
Hilbert space $L^2(\R)$ of square-integrable functions. In 1950, \textbf{Laurent Schwartz} developed
the theory of \emph{tempered distributions}, proving that generalized functions like Dirac's
$\delta(\omega)$ permit pure sinusoids, DC constants, and unit steps to possess rigorous,
well-defined Fourier transforms. The electrical engineering practice of writing $X(\jj\omega)$
highlights that the Fourier transform is simply the Laplace transform $X(s)$ evaluated along the
pure imaginary frequency axis $s = \jj\omega$.
\end{historybox}

\section{The continuous-time Fourier transform}
\label{ch07:ctft-defn}

\begin{definition}[Fourier transform pair]
\label{ch07:def-pair}
For a signal $x(t)$ the \emph{Fourier transform} (analysis equation) and its
\emph{inverse} (synthesis equation) are
\begin{equation}
\label{ch07:ftpair}
\boxed{\;X(\jj\omega) = \int_{-\infty}^{\infty} x(t)\, e^{-\jj\omega t}\, \dd t\;}
\qquad
\boxed{\;x(t) = \frac{1}{2\pi}\int_{-\infty}^{\infty} X(\jj\omega)\, e^{\jj\omega t}\, \dd\omega\;}
\end{equation}
We write $x(t) \overset{\FT}{\longleftrightarrow} X(\jj\omega)$, or $X = \FT\{x\}$ and
$x = \FT^{-1}\{X\}$.
\end{definition}

The two formulas are nearly mirror images. The forward transform multiplies by
$e^{-\jj\omega t}$ and integrates over time; the inverse multiplies by $e^{+\jj\omega t}$
and integrates over frequency, with a $1/(2\pi)$ out front. That lone factor of $2\pi$
is the only asymmetry, and it is the price of measuring frequency in radians per second.
(Some authors measure in cycles per second, $f = \omega/2\pi$; then both integrals carry
$e^{\mp \jj 2\pi f t}$ and the $2\pi$ disappears. We will note this convention where it
matters.)

\begin{whybox}
Why does $\int x(t)\,e^{-\jj\omega t}\,\dd t$ measure ``how much frequency $\omega$ is in
$x$''? Because $e^{-\jj\omega t}$ is a unit-length complex phasor spinning at rate
$-\omega$. Multiplying $x(t)$ by it \emph{de-rotates} any component of $x$ that spins at
$+\omega$, freezing it into a constant; every other component keeps spinning. Integrating
then averages: the frozen component survives and accumulates, while the still-spinning
components average to nearly zero over the whole line. The integral is a
\emph{correlation} of $x$ with a test phasor -- a projection of the signal onto the
sinusoid of frequency $\omega$. It is exactly the inner product
$\inner{x}{e^{\jj\omega t}}$ from the linear-algebra view of signals, now taken over the
whole real line.
\end{whybox}

\begin{happeningbox}
The inverse transform \eqref{ch07:ift} is a recipe: to rebuild $x(t)$, take a phasor
$e^{\jj\omega t}$ at \emph{every} frequency, weight it by the complex number
$X(\jj\omega)\,\dd\omega/(2\pi)$, and add them all up. A transient signal is therefore
an infinite blend of eternal, never-decaying sinusoids whose phases are arranged so that
they interfere constructively in one small region of time and destructively everywhere
else. What localizes a pulse in time is the \emph{phase} alignment of infinitely
many sine waves.
\end{happeningbox}

\paragraph{When does the integral exist?}
A clean sufficient condition is absolute integrability:
\begin{equation}
\int_{-\infty}^{\infty} \abs{x(t)}\, \dd t < \infty
\quad\Longrightarrow\quad X(\jj\omega) \text{ exists, bounded and continuous.}
\end{equation}
This covers every finite-energy pulse you will actually sample from a sensor. Many
useful idealizations -- a pure cosine, a constant, a step -- are \emph{not} absolutely
integrable, yet they have perfectly usable transforms once we allow Dirac deltas. We
treat those as limits and as distributions; the engineering benefit is large and the
bookkeeping is light.

\begin{connectionbox}
This is the same projection idea as the inner product in the linear-algebra companion,
pushed to a continuous ``basis.'' There, a vector $\vec v$ is reconstructed as
$\sum_k \inner{\vec v}{\vec e_k}\vec e_k$ over an orthonormal basis $\{\vec e_k\}$. Here
the ``basis vectors'' are the phasors $e^{\jj\omega t}$, indexed by a continuous label
$\omega$, and the sum becomes the integral \eqref{ch07:ift}. The catch is that
$e^{\jj\omega t}$ has infinite energy, so it is not a true element of the space $L^2$;
making the analogy airtight needs the theory of tempered distributions. The geometry,
though, is exactly the finite-dimensional one you already know.
\end{connectionbox}

\section{Magnitude and phase: reading a spectrum}
\label{ch07:mag-phase}

$X(\jj\omega)$ is complex, so at each frequency it carries two numbers. Write it in polar
form:
\begin{equation}
X(\jj\omega) = \abs{X(\jj\omega)}\, e^{\jj\,\angle X(\jj\omega)} .
\end{equation}
The \emph{magnitude spectrum} $\abs{X(\jj\omega)}$ says how strongly frequency $\omega$ is
present. The \emph{phase spectrum} $\angle X(\jj\omega)$ says \emph{where} that sinusoid
sits in time -- its alignment. Magnitude tells you the ingredients; phase tells you when
each ingredient arrives.

\begin{trapbox}
Beginners stare only at the magnitude spectrum and ignore phase. For \emph{recognizing}
a steady tone, magnitude is often enough. But phase is what localizes events in time. Two
signals can share an identical magnitude spectrum yet sound or look completely different:
a sharp drum hit and a smeared whoosh can have the same magnitude spectrum and opposite
phase behavior. If you ever scramble the phases of a recording while keeping magnitudes,
you destroy all transients and are left with noise-like mush, so phase cannot be ignored.
\end{trapbox}

\paragraph{Real signals have conjugate-symmetric spectra.}
If $x(t)$ is real (every physical sensor reading is), then
\begin{equation}
\label{ch07:herm}
X(-\jj\omega) = \overline{X(\jj\omega)} .
\end{equation}
This single fact has two consequences you will use constantly: the magnitude spectrum is
\emph{even} ($\abs{X(-\jj\omega)} = \abs{X(\jj\omega)}$) and the phase spectrum is
\emph{odd}. So the negative-frequency half of a real signal's spectrum is just a mirror
of the positive half -- it carries no new information, only the bookkeeping needed to
keep the synthesis real. This is why, in practice, we plot only $\omega \ge 0$.

\begin{figure}[H]
\centering
\begin{tikzpicture}
\begin{axis}[name=mag,width=6.7cm,height=4.2cm,
  xlabel={$\omega$},ylabel={$|X(\jj\omega)|$},
  xmin=-8,xmax=8,ymin=0,ymax=1.2,axis lines=left,ytick=\empty,
  title style={font=\small}, title={Magnitude (even)}]
\addplot[specred,thick,domain=-8:8,samples=160]{exp(-x^2/8)};
\end{axis}
\begin{axis}[name=ph,width=6.7cm,height=4.2cm,at={(mag.right of east)},anchor=left of west,xshift=1.2cm,
  xlabel={$\omega$},ylabel={$\angle X(\jj\omega)$},
  xmin=-8,xmax=8,ymin=-1.7,ymax=1.7,axis lines=middle,ytick=\empty,
  title style={font=\small}, title={Phase (odd)}]
\addplot[sigblue,thick,domain=-8:8,samples=160]{0.18*x};
\end{axis}
\end{tikzpicture}
\caption{For any real signal the magnitude spectrum is an even function and the phase
spectrum is odd. A linear phase like the one shown ($\angle X = -\omega t_0$ up to sign)
corresponds to a pure time shift.}
\label{ch07:fig-magphase}
\end{figure}

\begin{intuitionbox}
A \emph{linear} phase, $\angle X(\jj\omega) = -\omega t_0$, means ``every frequency is
delayed by the same amount $t_0$.'' Geometrically, delaying a signal in time rotates each
phasor in proportion to its frequency: fast phasors must rotate more to keep step. A
delay therefore shows up as a phase ramp, never as a change in magnitude. We prove this
as the time-shift property in Chapter~11, but the picture is worth keeping
now: \emph{shifting in time tilts the phase, it never touches the magnitude.}
\end{intuitionbox}

\section{A worked transform: the rectangle and the sinc}
\label{ch07:rect-sinc}

The most important transform pair in the subject is the rectangle and the
sinc. Once you understand it, much of sampling theory, filter design, and windowing will
follow naturally.

Let the unit-height rectangular pulse of total width $\tau$ be
\begin{equation}
x(t) = \rect\!\left(\frac{t}{\tau}\right) =
\begin{cases} 1, & \abs{t} < \tau/2,\\[2pt] 0, & \abs{t} > \tau/2. \end{cases}
\end{equation}
Its transform is a direct integral:
\begin{align}
X(\jj\omega)
&= \int_{-\tau/2}^{\tau/2} e^{-\jj\omega t}\, \dd t
= \left[\frac{e^{-\jj\omega t}}{-\jj\omega}\right]_{-\tau/2}^{\tau/2}
= \frac{e^{\jj\omega\tau/2} - e^{-\jj\omega\tau/2}}{\jj\omega}
\notag\\
&= \frac{2\sin(\omega\tau/2)}{\omega}
= \tau\,\frac{\sin(\omega\tau/2)}{\omega\tau/2}
= \tau\,\sinc\!\left(\frac{\omega\tau}{2}\right),
\label{ch07:rect-result}
\end{align}
using the unnormalized convention $\sinc(u) = \sin(u)/u$ with $\sinc(0)=1$. So
\begin{equation}
\boxed{\;\rect\!\left(\frac{t}{\tau}\right) \overset{\FT}{\longleftrightarrow}
\tau\,\sinc\!\left(\frac{\omega\tau}{2}\right).\;}
\end{equation}

\begin{figure}[H]
\centering
\begin{tikzpicture}
\begin{axis}[name=rt,width=6.4cm,height=4.4cm,
  xlabel={$t$},ylabel={$x(t)$},
  xmin=-2.4,xmax=2.4,ymin=-0.25,ymax=1.35,axis lines=middle,
  xtick={-1,1},xticklabels={$-\tau/2$,$\tau/2$},ytick=\empty,
  title style={font=\small}, title={Rectangle, width $\tau$}]
\addplot[sigblue,thick] coordinates {(-2.4,0)(-1,0)(-1,1)(1,1)(1,0)(2.4,0)};
\addplot[fillsig,opacity=0.35] coordinates {(-1,0)(-1,1)(1,1)(1,0)} \closedcycle;
\node[anchor=south east,font=\footnotesize] at (axis cs:0,1.02) {$1$};
\end{axis}
\begin{axis}[name=sc,width=7.2cm,height=4.4cm,at={(rt.right of east)},anchor=left of west,xshift=1.2cm,
  xlabel={$\omega$},ylabel={$X(\jj\omega)$},
  xlabel style={at={(axis cs:19,0)}, anchor=west},
  ylabel style={at={(axis cs:0,2.45)}, anchor=south},
  xmin=-19,xmax=19,ymin=-0.85,ymax=2.45,axis lines=middle,
  xtick={-3.1416,3.1416},
  xticklabels={$-\frac{2\pi}{\tau}$,$\frac{2\pi}{\tau}$},
  xticklabel style={font=\footnotesize, anchor=north, yshift=-3mm},
  ytick=\empty,
  title style={font=\small}, title={Sinc, first null at $2\pi/\tau$}]
\addplot[specred,thick,domain=-19:19,samples=199]{2*sin(deg(x+0.0001))/(x+0.0001)};
\node[anchor=south east,font=\footnotesize] at (axis cs:0,2.02) {$\tau$};
\end{axis}
\end{tikzpicture}
\caption{The fundamental pair. A rectangle of width $\tau$ (here $\tau=2$) transforms to
$\tau\,\sinc(\omega\tau/2)$. The first zero crossings sit at $\omega = \pm 2\pi/\tau$,
so a \emph{narrower} pulse gives a \emph{wider} sinc. The central peak has height $\tau$,
the area of the pulse.}
\label{ch07:fig-rectsinc}
\end{figure}

Two features of \eqref{ch07:rect-result} deserve a moment.

\begin{itemize}
  \item \textbf{The DC value is the area.} At $\omega=0$, $X(\jj 0) = \tau = \int x\,\dd t$.
        This is general: $X(\jj 0)$ always equals the total area of the signal. Frequency
        zero asks ``what is the net DC content?'' and the answer is the integral.
  \item \textbf{Nulls at multiples of $2\pi/\tau$.} The sinc crosses zero whenever
        $\omega\tau/2 = m\pi$, i.e. $\omega = 2\pi m/\tau$. A short pulse pushes those
        nulls outward: its energy is smeared across a broad band. A long pulse pulls them
        inward: its energy is concentrated near DC. This reciprocity leads to the
        uncertainty principle in Section~\ref{ch07:uncertainty}.
\end{itemize}

\begin{example}[A $20\,$ms gate]
\label{ch07:ex-gate}
Suppose a gesture detector marks a window of activity as a rectangular gate of width
$\tau = 20\,\text{ms} = 0.02\,\text{s}$. The first spectral null of that gate sits at
$\omega = 2\pi/\tau = 2\pi/0.02 \approx 314\,\text{rad/s}$, i.e. $f = 1/\tau = 50\,$Hz.
So merely \emph{turning a measurement on and off} over $20\,$ms injects spectral content
spread out to roughly $50\,$Hz and beyond -- a fact we must respect when we later filter
the gated signal, lest the gate's own spectrum masquerade as real motion.
\end{example}

\begin{whybox}
Why a sinc and not some other curve? Because a rectangle is the indicator of an interval,
and the transform of an indicator is the integral of a complex phasor over that interval.
Integrating $e^{-\jj\omega t}$ over a symmetric interval gives a difference of two phasors
$e^{\pm\jj\omega\tau/2}$, and the difference of two unit phasors equidistant from the real
axis is purely imaginary and equal to $2\jj\sin(\cdot)$. Divide by the $\jj\omega$ from
the integration and you are left with $2\sin(\omega\tau/2)/\omega$. The sinc is the
fingerprint of \emph{sharp edges}: any hard on/off boundary in time produces sinc-like
ripples in frequency.
\end{whybox}

\paragraph{Duality: it works the other way too.}
Because the analysis and synthesis integrals differ only by a sign and a $2\pi$, the pair
runs in reverse. A \emph{sinc in time} transforms to a \emph{rectangle in frequency}:
\begin{equation}
\label{ch07:dual}
\frac{\sin(\Omega t)}{\pi t} \overset{\FT}{\longleftrightarrow}
\rect\!\left(\frac{\omega}{2\Omega}\right)
=\begin{cases}1,&\abs{\omega}<\Omega,\\0,&\abs{\omega}>\Omega.\end{cases}
\end{equation}
This is the ideal low-pass filter: a brick wall in frequency that passes everything below
$\Omega$ and blocks everything above. Its time-domain shape is an infinitely long,
two-sided sinc -- which is why a \emph{perfect} brick-wall filter is non-causal and must
be approximated in practice. We will build that story in the filtering chapters; for now,
notice that duality hands us a second pair for free.

\begin{aestheticbox}
The rect/sinc duality captures a fundamental symmetry of the transform. The crispest possible thing
in time (a hard-edged box) is the most spread-out thing in frequency (an endless ripple),
and vice versa. The Fourier transform turns sharpness into spread and spread into
sharpness. Once you understand this, you can predict the
qualitative shape of a spectrum just by looking at how abrupt a signal is.
\end{aestheticbox}

\section{More essential transforms}
\label{ch07:more-pairs}

A small dictionary of pairs covers most of practice. Each one has a geometric interpretation.

\subsection{The impulse and the constant}
The Dirac delta $\delta(t)$ is the idealized instantaneous spike with unit area. By the
sifting property,
\begin{equation}
\FT\{\delta(t)\} = \int_{-\infty}^{\infty}\delta(t)\,e^{-\jj\omega t}\,\dd t = e^{-\jj\omega\cdot 0} = 1 .
\end{equation}
An impulse contains \emph{every} frequency equally and in phase -- a perfectly flat,
white spectrum. This is why striking a system with a sharp impulse and listening to the
response reveals the system's behavior at all frequencies at once (the impulse response,
central to LTI systems). Running the pair backward by duality gives the transform of a
constant:
\begin{equation}
\FT\{1\} = 2\pi\,\delta(\omega) .
\end{equation}
A constant (pure DC) is all area at frequency zero and nothing elsewhere -- an impulse in
the frequency domain.

\subsection{The complex exponential and the cosine}
Shifting the constant's spectrum sideways gives the eternal phasor:
\begin{equation}
\FT\{e^{\jj\omega_0 t}\} = 2\pi\,\delta(\omega-\omega_0) .
\end{equation}
A single never-ending sinusoid is a single spectral line -- an impulse at its own
frequency. Real cosines and sines are sums of two phasors, so they give two lines:
\begin{equation}
\label{ch07:cos}
\FT\{\cos(\omega_0 t)\} = \pi\,\delta(\omega-\omega_0) + \pi\,\delta(\omega+\omega_0),
\qquad
\FT\{\sin(\omega_0 t)\} = \frac{\pi}{\jj}\big[\delta(\omega-\omega_0) - \delta(\omega+\omega_0)\big].
\end{equation}
The pair of impulses at $\pm\omega_0$ is exactly the conjugate symmetry of
\eqref{ch07:herm} for a real signal. Cosine's two impulses are real and equal; sine's
are imaginary and opposite, encoding the $90^\circ$ phase difference between them
(Figure~\ref{ch07:fig-cosline}).

\begin{figure}[H]
\centering
\begin{tikzpicture}
\begin{axis}[width=11cm,height=4.4cm,
  xlabel={$\omega$},ylabel={$X(\jj\omega)$},
  xmin=-5,xmax=5,ymin=0,ymax=4.2,axis lines=middle,
  xtick={-2.5,2.5},xticklabels={$-\omega_0$,$\omega_0$},
  ytick={3.1416},yticklabels={$\pi$},
  title style={font=\small}, title={Spectrum of $\cos(\omega_0 t)$: two impulses}]
\addplot[specred,thick,mark=*,mark size=2pt,ycomb] coordinates {(-2.5,3.1416)(2.5,3.1416)};
\end{axis}
\end{tikzpicture}
\caption{A never-ending cosine has a line spectrum: two impulses of weight $\pi$ at
$\omega=\pm\omega_0$ (arrows drawn as stems, since a true impulse has infinite height and
finite area). Sine differs only in that its two impulses are imaginary and of opposite
sign.}
\label{ch07:fig-cosline}
\end{figure}

\subsection{The one-sided decaying exponential}
A causal decay $x(t) = e^{-at}\,u(t)$ with $a>0$ (here $u$ is the unit step) is the
prototype of a first-order system settling toward rest:
\begin{equation}
X(\jj\omega) = \int_0^{\infty} e^{-at}e^{-\jj\omega t}\,\dd t
= \frac{1}{a+\jj\omega},
\qquad
\abs{X(\jj\omega)} = \frac{1}{\sqrt{a^2+\omega^2}} .
\end{equation}
The magnitude falls off as $1/\omega$ at high frequency -- a gentle low-pass shape. The
parameter $a$ sets the \emph{corner}: at $\omega=a$ the magnitude has dropped by
$1/\sqrt2$. Fast decay (large $a$, short time constant) means a wide spectrum; slow decay
means a narrow one. The same reciprocity appears again (Figure~\ref{ch07:fig-decay}).

\begin{figure}[H]
\centering
\begin{tikzpicture}
\begin{axis}[name=de1,width=6.4cm,height=4.2cm,
  xlabel={$t$},ylabel={$x(t)$},xmin=-0.6,xmax=4,ymin=0,ymax=1.15,
  axis lines=left,ytick={1},title style={font=\small},title={Causal decay $e^{-at}u(t)$}]
\addplot[sigblue,thick,domain=0:4,samples=120]{exp(-1.5*x)};
\addplot[fillsig,opacity=0.25,domain=0:4,samples=120]{exp(-1.5*x)}\closedcycle;
\end{axis}
\begin{axis}[name=de2,width=6.6cm,height=4.2cm,at={(de1.east)},anchor=west,xshift=1.2cm,
  xlabel={$\omega$},ylabel={$|X(\jj\omega)|$},xmin=0,xmax=8,ymin=0,ymax=0.75,
  axis lines=left,xtick={1.5},xticklabels={$a$},ytick={0.4714},
  yticklabels={$\tfrac{1}{\sqrt2}\,\tfrac1a$},title style={font=\small},title={Magnitude, corner at $\omega=a$}]
\addplot[specred,thick,domain=0:8,samples=160]{1/sqrt(1.5^2+x^2)};
\addplot[gray,dashed] coordinates {(1.5,0)(1.5,0.4714)};
\addplot[gray,dashed] coordinates {(0,0.4714)(1.5,0.4714)};
\end{axis}
\end{tikzpicture}
\caption{The one-sided decaying exponential and its magnitude spectrum
$1/\sqrt{a^2+\omega^2}$. The response is flat near DC (height $1/a$) and rolls off as
$1/\omega$; at the corner $\omega=a$ it has fallen to $1/\sqrt2$ of its peak. A faster
decay (larger $a$) pushes the corner higher and widens the band.}
\label{ch07:fig-decay}
\end{figure}

\subsection{The Gaussian: the shape that transforms to itself}
The Gaussian pulse $x(t) = e^{-t^2/(2\sigma^2)}$ transforms to another Gaussian:
\begin{equation}
\label{ch07:gauss}
e^{-t^2/(2\sigma^2)} \overset{\FT}{\longleftrightarrow}
\sigma\sqrt{2\pi}\; e^{-\sigma^2\omega^2/2}.
\end{equation}
A wide time-Gaussian (large $\sigma$) becomes a narrow frequency-Gaussian (width
$1/\sigma$), and conversely. No edges, no ripples -- the Gaussian is the smoothest
possible bump and pays for that smoothness with the smoothest possible spectrum. It is
also the unique shape that \emph{simultaneously minimizes} the time-frequency spread, as
we will see next, which is why Gaussian windows are prized in spectral analysis.

\begin{figure}[H]
\centering
\begin{tikzpicture}
\begin{axis}[name=g1,width=6.4cm,height=4.2cm,
  xlabel={$t$},ylabel={$x(t)$},xmin=-4,xmax=4,ymin=0,ymax=1.15,
  axis lines=left,ytick=\empty,title style={font=\small},title={Narrow Gaussian (small $\sigma$)}]
\addplot[sigblue,thick,domain=-4:4,samples=120]{exp(-x^2/0.5)};
\addplot[fillsig,opacity=0.25,domain=-4:4,samples=120]{exp(-x^2/0.5)}\closedcycle;
\end{axis}
\begin{axis}[name=g2,width=6.4cm,height=4.2cm,at={(g1.right of east)},anchor=left of west,xshift=1.2cm,
  xlabel={$\omega$},ylabel={$|X(\jj\omega)|$},xmin=-8,xmax=8,ymin=0,ymax=1.15,
  axis lines=left,ytick=\empty,title style={font=\small},title={Wide Gaussian spectrum}]
\addplot[specred,thick,domain=-8:8,samples=120]{exp(-x^2/8)};
\addplot[specred!30,opacity=0.4,domain=-8:8,samples=120]{exp(-x^2/8)}\closedcycle;
\end{axis}
\end{tikzpicture}
\caption{A Gaussian transforms to a Gaussian. Squeeze it in time and it spreads in
frequency by the reciprocal amount; their widths multiply to a constant.}
\label{ch07:fig-gauss}
\end{figure}

\begin{table}[H]
\centering
\caption{A working dictionary of Fourier transform pairs (radian-frequency convention).}
\label{ch07:tbl-pairs}
\begin{tabular}{@{}lll@{}}
\toprule
Name & $x(t)$ & $X(\jj\omega)$ \\
\midrule
Impulse        & $\delta(t)$                 & $1$ \\
Constant       & $1$                         & $2\pi\,\delta(\omega)$ \\
Phasor         & $e^{\jj\omega_0 t}$         & $2\pi\,\delta(\omega-\omega_0)$ \\
Cosine         & $\cos(\omega_0 t)$          & $\pi[\delta(\omega-\omega_0)+\delta(\omega+\omega_0)]$ \\
Rectangle      & $\rect(t/\tau)$             & $\tau\,\sinc(\omega\tau/2)$ \\
Sinc           & $\dfrac{\sin(\Omega t)}{\pi t}$ & $\rect(\omega/2\Omega)$ \\
Causal decay   & $e^{-at}u(t),\ a>0$         & $\dfrac{1}{a+\jj\omega}$ \\
Gaussian       & $e^{-t^2/(2\sigma^2)}$      & $\sigma\sqrt{2\pi}\,e^{-\sigma^2\omega^2/2}$ \\
\bottomrule
\end{tabular}
\end{table}

\section{The uncertainty principle: time and frequency cannot both be sharp}
\label{ch07:uncertainty}

Every pair in Table~\ref{ch07:tbl-pairs} exhibits the same reciprocal law. Narrow in time means broad in frequency.
The rectangle's sinc nulls march outward as $\tau$ shrinks; the decaying exponential's
corner climbs as the decay quickens; the Gaussian's frequency width is exactly $1/\sigma$.
The pattern is a theorem, not a coincidence of these examples.

Define the effective \emph{duration} and \emph{bandwidth} of a finite-energy signal as
the standard deviations of its energy density in each domain. With energy
$E = \int\abs{x(t)}^2\dd t = \frac{1}{2\pi}\int\abs{X(\jj\omega)}^2\dd\omega$ (Parseval,
Section~\ref{ch07:parseval}), and assuming the signal centered at $t=0$ and $\omega=0$,
\begin{equation}
(\Delta t)^2 = \frac{1}{E}\int t^2\abs{x(t)}^2\,\dd t,
\qquad
(\Delta\omega)^2 = \frac{1}{2\pi E}\int \omega^2\abs{X(\jj\omega)}^2\,\dd\omega .
\end{equation}

\begin{theorem}[Time-bandwidth uncertainty]
\label{ch07:thm-uncertainty}
For any finite-energy signal,
\begin{equation}
\Delta t \,\cdot\, \Delta\omega \;\ge\; \tfrac{1}{2},
\end{equation}
with equality \emph{only} for a Gaussian pulse.
\end{theorem}

The proof is a Cauchy--Schwarz argument applied to $t\,x(t)$ and $x'(t)$, using
integration by parts and the differentiation property $\FT\{x'(t)\} = \jj\omega\,X(\jj\omega)$;
we sketch it in the exercises. The takeaway needs no algebra.

\begin{intuitionbox}
To pin down \emph{when} something happens you must look at a short slice of time -- but a
short slice cannot contain many cycles of any single sinusoid, so its frequency is
blurry. To pin down \emph{what frequency} a tone is, you must watch it for many cycles --
but watching for a long time blurs out exactly \emph{when} it occurred. Time precision
and frequency precision pull in opposite directions. You can trade one for the other, but their
product can never beat the floor $\tfrac12$.
\end{intuitionbox}

\begin{figure}[H]
\centering
\begin{tikzpicture}
\begin{axis}[name=a,width=6.4cm,height=3.9cm,xlabel={$t$},ylabel={},
  xmin=-4,xmax=4,ymin=0,ymax=1.1,axis lines=left,ytick=\empty,
  title style={font=\small}, title={Short pulse}]
\addplot[sigblue,thick,domain=-4:4,samples=120]{exp(-x^2/0.25)};
\end{axis}
\begin{axis}[name=b,width=6.4cm,height=3.9cm,at={(a.right of east)},anchor=left of west,xshift=1.1cm,
  xlabel={$\omega$},ylabel={},xmin=-12,xmax=12,ymin=0,ymax=1.1,axis lines=left,ytick=\empty,
  title style={font=\small}, title={$\Rightarrow$ broad spectrum}]
\addplot[specred,thick,domain=-12:12,samples=120]{exp(-x^2/16)};
\end{axis}
\end{tikzpicture}

\vspace{0.4em}

\begin{tikzpicture}
\begin{axis}[name=c,width=6.4cm,height=3.9cm,xlabel={$t$},ylabel={},
  xmin=-4,xmax=4,ymin=0,ymax=1.1,axis lines=left,ytick=\empty,
  title style={font=\small}, title={Long pulse}]
\addplot[sigblue,thick,domain=-4:4,samples=120]{exp(-x^2/4)};
\end{axis}
\begin{axis}[name=d,width=6.4cm,height=3.9cm,at={(c.right of east)},anchor=left of west,xshift=1.1cm,
  xlabel={$\omega$},ylabel={},xmin=-12,xmax=12,ymin=0,ymax=1.1,axis lines=left,ytick=\empty,
  title style={font=\small}, title={$\Rightarrow$ narrow spectrum}]
\addplot[specred,thick,domain=-12:12,samples=120]{exp(-x^2/1)};
\end{axis}
\end{tikzpicture}
\caption{The uncertainty tradeoff in pictures. Squeezing the pulse in time (top) forces
the spectrum to spread; stretching it (bottom) sharpens the spectrum. Their widths cannot
both be small.}
\label{ch07:fig-uncertainty}
\end{figure}

\begin{applicationbox}
This is the governing constraint of every spectrogram. To see a fast event you use a
short analysis window and accept coarse frequency resolution; to resolve two close tones
you use a long window and accept smeared timing. There is no setting that gives both. The
short-time Fourier transform chapter is largely about choosing the
best compromise on this tradeoff for the job at hand.
\end{applicationbox}

\begin{goalbox}
The uncertainty principle is the first hard physical limit on our wrist controller, and
it is worth stating bluntly. Suppose we want the instrument to respond to a flick of the
wrist within $20\,$ms so it feels instantaneous to a player. By
Theorem~\ref{ch07:thm-uncertainty}, a $20\,$ms analysis window can resolve frequency only
to roughly $\Delta f \sim 1/(2\pi\cdot 0.02) \approx 8\,$Hz at the theoretical floor -- in
practice, with a real tapering window, more like $50$--$80\,$Hz. That is plenty to tell a
sharp tap from a slow sweep, but it is \emph{not} enough to separate two motions that
differ only by a few Hz in their tremor spectrum. Low latency and fine frequency
discrimination pull in opposite directions; the controller's gesture set must be designed
to live within that budget rather than fight it. We return to this when we lay out the
capstone.
\end{goalbox}

\section{Parseval and the conservation of energy}
\label{ch07:parseval}

Beyond reorganizing a signal's information across frequency, the transform preserves its
\emph{energy}. This is the continuous analogue of the fact that a rotation does not change
the length of a vector -- the Fourier transform is (up to the $2\pi$) an isometry.

\begin{theorem}[Parseval / Rayleigh energy theorem]
\label{ch07:thm-parseval}
For finite-energy signals,
\begin{equation}
\int_{-\infty}^{\infty}\abs{x(t)}^2\,\dd t
= \frac{1}{2\pi}\int_{-\infty}^{\infty}\abs{X(\jj\omega)}^2\,\dd\omega .
\end{equation}
More generally the inner product is preserved:
$\int x(t)\overline{y(t)}\,\dd t = \frac{1}{2\pi}\int X(\jj\omega)\overline{Y(\jj\omega)}\,\dd\omega$.
\end{theorem}

\begin{whybox}
The general identity follows by writing one signal through the inverse transform and
swapping the order of integration:
\[
\int x(t)\overline{y(t)}\,\dd t
= \int x(t)\left[\frac{1}{2\pi}\int \overline{Y(\jj\omega)}\,e^{-\jj\omega t}\,\dd\omega\right]\dd t
= \frac{1}{2\pi}\int \overline{Y(\jj\omega)}\left[\int x(t)e^{-\jj\omega t}\dd t\right]\dd\omega,
\]
and the inner bracket is exactly $X(\jj\omega)$. Setting $y=x$ gives the energy form. In
the language of the linear-algebra companion, the transform preserves inner products, so
it preserves lengths and angles: it is a unitary map (after rescaling by $1/\sqrt{2\pi}$).
\end{whybox}

The quantity $\abs{X(\jj\omega)}^2$ is the \emph{energy spectral density}: integrate it
over a band of frequencies and you get the energy the signal carries in that band. This
is how you answer questions like ``what fraction of this twitch's energy lives below
$30\,$Hz?'' -- you integrate the density there and divide by the total.

\begin{example}[Energy of a rectangular pulse, both ways]
\label{ch07:ex-parseval}
The pulse $\rect(t/\tau)$ has time-domain energy $\int_{-\tau/2}^{\tau/2}1^2\,\dd t = \tau$.
Parseval says the frequency side agrees:
\[
\frac{1}{2\pi}\int_{-\infty}^{\infty}\left[\tau\,\sinc\!\Big(\tfrac{\omega\tau}{2}\Big)\right]^2\dd\omega
= \frac{\tau^2}{2\pi}\cdot\frac{2\pi}{\tau} = \tau,
\]
using the standard integral $\int_{-\infty}^{\infty}\sinc^2(\omega\tau/2)\,\dd\omega = 2\pi/\tau$.
The numbers match. The energy of a pulse is the same whether you count it in time or sum
the squared magnitudes of its frequency content.
\end{example}

\begin{connectionbox}
Parseval is the integral cousin of the Pythagorean theorem. In a finite-dimensional inner
product space, $\norm{\vec v}^2 = \sum_k\abs{\inner{\vec v}{\vec e_k}}^2$ over an
orthonormal basis. Here the ``coordinates'' of $x$ are its transform values
$X(\jj\omega)$, the sum becomes an integral, and the $1/(2\pi)$ is the normalization that
makes the phasor ``basis'' behave like an orthonormal one. Energy in time equals energy
in frequency for the same reason that a vector's length does not depend on which
orthonormal axes you measure it along.
\end{connectionbox}

\section{A computational sanity check}
\label{ch07:numeric}

We have not yet built the discrete machinery (that is the next chapter), but it is worth
seeing that the rect/sinc pair is not an abstraction. The following snippet samples a
rectangle densely, takes a crude numerical transform, and confirms the sinc shape and the
first-null location.

\begin{lstlisting}[language=Python,caption={Numerically confirming rect to sinc.\label{ch07:lst-rect}}]
import numpy as np

fs   = 2000.0            # samples per second
t    = np.arange(-1.0, 1.0, 1/fs)
tau  = 0.1               # 100 ms rectangle
x    = (np.abs(t) < tau/2).astype(float)

# crude continuous-transform estimate: Riemann sum of x(t) e^{-j w t} dt
w    = 2*np.pi*np.linspace(-200, 200, 4001)        # rad/s grid
X    = (x[None, :] * np.exp(-1j*np.outer(w, t))).sum(axis=1) / fs

print("first null (predicted) =", 2*np.pi/tau, "rad/s")   # ~ 62.8
print("peak |X| vs tau:", np.abs(X).max(), "vs", tau)     # both ~ 0.1
\end{lstlisting}

The printed peak magnitude lands on $\tau = 0.1$ (the pulse area, as predicted by the DC
rule), and the first null appears at $\omega = 2\pi/\tau \approx 62.8\,$rad/s. The Riemann
sum is exactly the analysis integral \eqref{ch07:ftpair} discretized; the next chapter
makes this precise with the discrete Fourier transform and the FFT.

\begin{buildbox}
On a microcontroller you will never evaluate the continuous integral; you will sample the
sensor at some rate $f_s$ and run a discrete transform. But the continuous picture is
still your design guide. Before writing a line of firmware, sketch the \emph{expected}
continuous spectrum of the motion you care about: a wrist flick is roughly a single bump
a few tens of milliseconds wide, so by the rect/sinc and Gaussian reasoning its energy
lives mostly below $\sim 20\,$Hz. That sketch tells you that a sampling rate of, say,
$200\,$Hz is comfortably more than enough for the gross motion, and it tells you where to
put the corner of your anti-alias and smoothing filters.
\end{buildbox}

\begin{goalbox}
A word about what the spectrum can and cannot tell us. A single wrist-mounted IMU
(accelerometer plus gyroscope) measures the gross motion of the forearm and hand as a
rigid-ish body. Its Fourier spectrum cleanly separates a deliberate flick from hand
tremor (a few Hz) from electrical hum, and that is genuinely useful for triggering
discrete gestures. What it \emph{cannot} do, no matter how cleverly you transform it, is
recover the angle of an individual finger joint: those motions are largely invisible to a
sensor strapped to the wrist. Continuous finger pose needs sensors \emph{on} the fingers
(instrumented rings or gloves); finger-level intent without rings needs surface EMG or
force myography (FMG) on the forearm, which sense muscle activity rather than bulk motion.
The Fourier transform is the right lens for all of these signals -- but it cannot conjure
information the sensor never captured. We keep this limit front and center as the build
progresses; the chapters on EMG/FMG and on gesture recognition spell out the precise
boundaries.
\end{goalbox}

\section{What to carry forward}
\label{ch07:recap}

\begin{recoverybox}
The Fourier transform is the Fourier series with the period sent to infinity: the
discrete harmonic ladder becomes a continuous spectrum $X(\jj\omega)$. The forward
transform projects $x(t)$ onto each eternal phasor $e^{\jj\omega t}$ (an inner product
over the whole line); the inverse reassembles $x$ from those projections. Magnitude says
how much of each frequency is present, phase says when it arrives, and for real signals
the spectrum is conjugate-symmetric so only $\omega\ge 0$ carries new information. The
rect/sinc pair encodes the main lesson -- sharpness in one domain is spread in the
other -- which becomes the uncertainty principle $\Delta t\,\Delta\omega\ge\tfrac12$,
saturated only by the Gaussian. Parseval guarantees energy is conserved across the two
domains. Next we discretize all of this into the DFT, the version that actually runs on
hardware.
\end{recoverybox}

\begin{connectionbox}
\textbf{Capabilities, Downsides, and Forward Roadmap}
\begin{itemize}
  \item \textbf{What this made possible}: The continuous Fourier transform removes the requirement
  of periodicity, allowing us to represent transient, non-repeating continuous signals of finite
  energy. It provides continuous spectral densities $X(\jj\omega)$, establishes the reciprocal
  time-bandwidth uncertainty relation $\Delta t\,\Delta\omega \ge 1/2$, and via Plancherel's theorem,
  proves that total signal energy is conserved across time and frequency.
  \item \textbf{Limitations}: The continuous transform is purely a mathematical
  idealization that cannot run directly on a digital computer:
  \begin{enumerate}
    \item \emph{Continuous integrals}: Evaluating $\int_{-\infty}^\infty x(t)e^{-\jj\omega t}\,\dd t$
    requires knowing $x(t)$ at uncountably many continuous instants over an infinite duration,
    impossible for digital microprocessors with finite discrete memory.
    \item \emph{Global loss of timing}: The transform is global: $X(\jj\omega)$ integrates over the
    entire history of time from $-\infty$ to $+\infty$. If a high-frequency burst occurs at $t = 2.5\,$s,
    $X(\jj\omega)$ shows high-frequency energy, but gives no indication of \emph{when} it happened.
  \end{enumerate}
  \item \textbf{Where we go next}: Chapter 8 discretizes time and frequency simultaneously, producing
  the Discrete Fourier Transform (DFT) that operates on finite buffers of $N$ digital samples.
  Chapter 9 introduces the Cooley--Tukey FFT algorithm to evaluate the DFT in $\mathcal{O}(N\log N)$
  time; and Chapter 10 develops the Short-Time Fourier Transform (STFT) and spectrograms to recover
  time-localization.
\end{itemize}
\end{connectionbox}

\section{Exercises}
\label{ch07:exercises}

\begin{exercise}[Area and DC value]
\label{ch07:ex-dc}
Without computing a full transform, state $X(\jj 0)$ for (a) a triangular pulse of base
$2\tau$ and peak height $1$, and (b) $e^{-at}u(t)$ with $a>0$. In each case explain why
$X(\jj 0)$ equals the time-domain area.
\end{exercise}

\begin{exercise}[Scaling and the reciprocity of widths]
\label{ch07:ex-scale}
Using the rect/sinc pair, find the transform of $\rect(t/(2\tau))$ and of
$\rect(2t/\tau)$. Where is the first null in each? Confirm that doubling the pulse width
halves the null frequency, and halving the width doubles it.
\end{exercise}

\begin{exercise}[Time shift tilts the phase]
\label{ch07:ex-shift}
Starting from Definition~\ref{ch07:def-pair}, show that
$\FT\{x(t-t_0)\} = e^{-\jj\omega t_0}X(\jj\omega)$. Conclude that a pure delay leaves
$\abs{X(\jj\omega)}$ unchanged and adds the linear phase $-\omega t_0$. Relate this to
Figure~\ref{ch07:fig-magphase}.
\end{exercise}

\begin{exercise}[Cosine as two lines]
\label{ch07:ex-cos}
Write $\cos(\omega_0 t) = \tfrac12 e^{\jj\omega_0 t} + \tfrac12 e^{-\jj\omega_0 t}$ and use
$\FT\{e^{\jj\omega_0 t}\}=2\pi\delta(\omega-\omega_0)$ to derive \eqref{ch07:cos}. Then do
the same for $\sin(\omega_0 t)$ and explain why its impulses are imaginary while cosine's
are real.
\end{exercise}

\begin{exercise}[Decaying exponential, magnitude and phase]
\label{ch07:ex-exp}
For $x(t) = e^{-at}u(t)$, $a>0$, you found $X(\jj\omega) = 1/(a+\jj\omega)$. Sketch
$\abs{X(\jj\omega)}$ and $\angle X(\jj\omega)$ for $\omega\ge 0$. At what $\omega$ does
the magnitude fall to $1/\sqrt2$ of its DC value, and what is the phase there?
\end{exercise}

\begin{exercise}[Parseval as a tool]
\label{ch07:ex-parseval2}
Use Parseval's theorem to evaluate $\int_{-\infty}^{\infty}\dfrac{\sin^2(\Omega t)}{(\pi t)^2}\,\dd t$
by recognizing the integrand as $\abs{x(t)}^2$ for the ideal low-pass sinc of
\eqref{ch07:dual}. (Hint: its transform is a rectangle of height $1$ and width $2\Omega$.)
\end{exercise}

\begin{exercise}[The Gaussian saturates uncertainty]
\label{ch07:ex-gauss}
For $x(t)=e^{-t^2/(2\sigma^2)}$ with transform \eqref{ch07:gauss}, compute $\Delta t$ and
$\Delta\omega$ (the energy-density standard deviations). Show $\Delta t = \sigma/\sqrt2$,
$\Delta\omega = 1/(\sigma\sqrt2)$, and hence $\Delta t\,\Delta\omega = \tfrac12$ exactly.
\end{exercise}

\begin{exercise}[Sketch the uncertainty proof]
\label{ch07:ex-uncertainty}
Fill in the Cauchy--Schwarz argument for Theorem~\ref{ch07:thm-uncertainty}: apply the
inequality $\abs{\int t\,x(t)\,\overline{x'(t)}\,\dd t}^2 \le \big(\int t^2\abs{x}^2\big)\big(\int\abs{x'}^2\big)$,
use integration by parts on the left to get $-\tfrac12\int\abs{x}^2\,\dd t$, and use
$\int\abs{x'(t)}^2\dd t = \frac{1}{2\pi}\int\omega^2\abs{X(\jj\omega)}^2\dd\omega$ on the
right. Assemble these into $\Delta t\,\Delta\omega\ge\tfrac12$.
\end{exercise}

\begin{exercise}[Numerical experiment]
\label{ch07:ex-numeric}
Adapt Listing~\ref{ch07:lst-rect} to a Gaussian $x(t)=e^{-t^2/(2\sigma^2)}$. Verify
numerically that the transform is again a Gaussian and that as you halve $\sigma$ the
spectral width doubles. Plot the product of the measured time and frequency widths and
confirm it stays near the floor $\tfrac12$.
\end{exercise}

\begin{exercise}[Gating a real signal]
\label{ch07:ex-gate2}
A $50\,$Hz cosine is observed only through a rectangular window of width $\tau=40\,$ms.
Using the convolution-in-frequency idea (multiplication in time corresponds to
convolution of spectra), describe qualitatively what happens to the two impulses of
\eqref{ch07:cos}: where do they end up, and how wide does each become? What does this say
about how cleanly you can read a tone's frequency from a short observation? Connect your
answer to Example~\ref{ch07:ex-gate}.
\end{exercise}
